\section{Discussion}

\subsection{Compatibility with Gaussian Noise}

A key factor behind the success of MeanFlow in our experiments lies in the compatibility between Gaussian noise and the flow-based generative modeling framework. Gaussian noise is smooth, unimodal, and isotropic, ensuring that the corrupted data distribution remains mathematically well-behaved (like smooth and differentiable), as required by the continuity equation underlying Flow Matching and MeanFlow. In our experiments, noisy inputs are generated by $\text{noisy-img} = \text{clean-img} + \sigma \epsilon$, which corresponds to a pixel-wise affine transformation of independent one-dimensional Gaussians. This preserves the smoothness and structure of the distribution, guaranteeing that the induced velocity fields remain continuous and amenable to stable transport. Conversely, when we clip each pixel of the noisy images to $[0, 1]$ before training, the probability mass of pixels exceeding boundaries accumulates at two ends, creating sharp discontinuities and $\delta$-like densities at the boundaries. Such modification breaks the smoothness assumption, violates the continuity requirement, and finally leads to failure, illustrating the importance of a well-behaved data distribution for both Flow Matching and MeanFlow.

\subsection{Architectural Considerations}

Our adoption of the Diffusion Transformer (DiT) as the backbone network introduces additional insights. Unlike convolutional networks that emphasize local receptive fields, transformers leverage self-attention to capture global dependencies across the entire image. This ability is particularly beneficial when dealing with Gaussian noise, which corrupts all pixels uniformly. By allowing information from distant parts of the image to influence local denoising decisions, DiT helps preserve the overall digit structure while recovering fine details. Nevertheless, transformers incur higher computational costs compared to CNNs or UNet, raising questions about scalability. Future work should explore trade-offs between global context modeling and efficiency, potentially through hybrid CNN-transformer architectures.

\subsection{Limitations and Future Directions}

Despite promising results, the present study has several limitations for further investigation. First, experiments are restricted to MNIST, a relatively simple dataset. Extending to more complex datasets such as CIFAR-10 or Tiny-ImageNet would better demonstrate scalability. Second, only Gaussian noise is considered, whereas real-world noise often exhibits non-Gaussian or structured characteristics (e.g., sensor noise, occlusions). Finally, we have not conducted a systematic computational comparison between MeanFlow and other denoising methods. From the model perspective, the current architecture also presents challenges. While MeanFlow achieves efficient sampling, the JVP operation incurs significant computational overhead during training. Additionally, the fixed guidance scale in CFG somewhat limits its flexibility. Future work would address these limitation by exploring more computationally efficient and effective training target, adaptive guidance strategies, and testing the framework on larger and more diverse datasets with complex noise patterns. Such work will provide a more comprehensive understanding of MeanFlow’s abilities and applications.
